\chapter{Benchmark externo (zero-shot)}
\label{ch:benchmark}

Este capítulo corresponde al notebook \texttt{06\_benchmark\_externo}. El
modelo mBERT afinado en el corpus local se evalúa sin reentrenamiento sobre
datasets públicos en español/latam hospedados en Hugging Face
~\cite{socialtox,sinodio,contextualizedhate,textdetox,aliaes,exist}.

\section{Protocolo}

Para cada dataset se proyectan las etiquetas originales a un esquema
compatible (multiclase o binario según disponibilidad), se limita el tamaño
según la configuración y se calculan F1-macro, F1-binario y accuracies.

\section{Resultados}

La Tabla~\ref{tab:benchmark} y la Figura~\ref{fig:benchmark} resumen los seis
datasets evaluados con estado \texttt{ok}. El mejor F1-macro zero-shot es
$0{,}313$ (contextualized hate speech); el peor, $0{,}139$ (SinOdio-LATAM).
En métrica binaria algunos datasets alcanzan F1 $\approx 0{,}70$, pero el
desempeño multiclase permanece lejos del F1-macro local de $0{,}778$.

\begin{table}[H]
\centering
\caption{Benchmark zero-shot de mBERT-SV sobre datasets externos.}
\label{tab:benchmark}
\begin{tabular}{lcccc}
\toprule
Dataset & n & F1-macro & F1-binario & Accuracy \\
\midrule
contextualized\_hate\_speech & 1614 & 0.3127 & 0.7110 & 0.3501 \\
SocialTOX & 968 & 0.2658 & 0.5802 & 0.4143 \\
SinOdio-LATAM-Regional-HateSpeech & 2000 & 0.1394 & 0.3171 & 0.2690 \\
multilingual\_toxicity\_dataset & 2000 & 0.2693 & 0.7023 & 0.4450 \\
ALIA-es-discriminative-hate-speech & 3000 & 0.2409 & 0.3826 & 0.4263 \\
exist\_dataset\_es & 1500 & 0.1716 & 0.5655 & 0.2933 \\
\bottomrule
\end{tabular}

\end{table}

\begin{figure}[H]
\centering
\includegraphics[width=0.95\textwidth]{05_benchmark_externo.png}
\caption{F1-macro y F1-binario en el benchmark externo.}
\label{fig:benchmark}
\end{figure}

\section{Discusión}

La caída de desempeño confirma un fuerte \emph{domain shift}: el modelo está
especializado en el dialecto, la distribución de clases y las plataformas del
corpus salvadoreño. La transferencia zero-shot a odio contextualizado argentino,
toxicidad social española o datasets de sexismo (EXIST) no es suficiente para
uso productivo sin adaptación adicional (fine-tuning, few-shot o alineación de
etiquetas más fina).
